\subsection{Paired variations of the central passage}
\label{sec:paired-variation}

The two graphs start at the same central height. Pairing their variations
at equal \(z\)-levels cancels the quadratic boundary term from the
endpoint derivative ratio. We use the remaining terms to prove
contraction when \(q_1\le0\), and then at sufficiently large central
heights when \(q_1>0\).

Reduce the capture data of \cref{prop:isolation-capture-new} so that
\begin{equation}
\label{eq:small-central-height-new}
 0<h<h_0<\frac{z_*}{2}.
\end{equation}
This is possible because the central crossing of the organizing graphic is
$(\eta,z)=(0,0)$.

Let $L_p$ denote the local map from the positive endpoint window to the
negative endpoint window obtained by following the central passage.
Expressing both endpoints in terms of the same central height gives,
in the coordinate $I=z_*-\eta^2/2$,
\begin{equation}
\label{eq:local-map-derivative-new}
 L_p'(I_+(h))=\frac{I_-'(h)}{I_+'(h)}>0.
\end{equation}

\begin{lemma}[Paired-variation identity]
\label{lem:paired-variation-identity-new}
Let $U_\pm$ be the graph solutions in
\eqref{eq:physical-graph-equation-new}.  Define
\begin{equation}
\label{eq:J-sigma-new}
 J_\sigma(z)=B_\eta(\eta_\sigma,z,p)-\frac{q_1}{z}
 -\frac{\eta_\sigma^2}{z}
  \bigl(3R_q(\eta_\sigma,p)
       +\eta_\sigma\partial_\eta R_q(\eta_\sigma,p)\bigr).
\end{equation}
Then
\begin{equation}
\label{eq:paired-log-derivative-new}
 \log L_p'(I_+(h))
 =-\int_h^{z_*}
 \left(\frac{J_-(z)}{\sqrt{2U_-(z;h)}}
       +\frac{J_+(z)}{\sqrt{2U_+(z;h)}}\right)\dd z.
\end{equation}
The integral is improper only at $z=h$ and is convergent.
\end{lemma}

\begin{proof}
Differentiate the right-hand side of
\eqref{eq:physical-graph-equation-new} with respect to $U$.  Since
$\eta_\sigma=\sigma\sqrt{2U}$, direct differentiation of
\eqref{eq:q-expansion-new} gives
\begin{equation}
\label{eq:F-U-derivative-new}
 \partial_U\mathcal G_{\sigma,p}
 =-\frac{2q_2}{z}
 +\frac{\sigma}{\sqrt{2U}}\left\{
 B_\eta-\frac{q_1}{z}
 -\frac{\eta^2}{z}(3R_q+\eta R_{q,\eta})\right\}.
\end{equation}
We compare the endpoint variations using
\eqref{eq:moving-height-variation-new}. Their initial factor $f_p(h)/h$
is common to both signs and cancels in the quotient. The term $-2q_2/z$
in \eqref{eq:F-U-derivative-new} cancels at each common \(z\)-level.
Subtracting the two remaining integrands gives
\eqref{eq:paired-log-derivative-new}. As in
\eqref{eq:moving-height-variation-new}, the singularity at \(z=h\)
is integrable.
\end{proof}

\begin{proposition}[Signed contraction]
\label{prop:signed-contraction-new}
If $q_1\le0$, then
\begin{equation}
\label{eq:signed-local-contraction-new}
 0<L_p'<1
\end{equation}
on the whole central return interval. The return is strictly
contracting in the endpoint coordinate $I$, and $\Delta_p$ is strictly
decreasing in $h$. In particular, $\Delta_p$ has at most one zero.
\end{proposition}

\begin{proof}
The sign conditions \eqref{eq:normal-sign-margins-new} and $q_1\le0$ give
$J_\sigma\ge b_1$ along both halves. The right-hand side of
\eqref{eq:paired-log-derivative-new} is therefore strictly negative,
which proves \eqref{eq:signed-local-contraction-new}. Composing with
the exterior transition in the endpoint coordinate and using
\eqref{eq:exterior-uniform-contraction-new} gives
$0<(\Pi_p\circ L_p)'<1$. In terms of the central height, this is
\[
 \Delta_p'(h)=I_+'(h)
   \left\{\Pi_p'(I_-(h))\frac{I_-'(h)}{I_+'(h)}-1\right\}<0.
\]
Thus the displacement has at most one zero.
\end{proof}

For $q_1>0$, the term $-q_1/z$ in \eqref{eq:J-sigma-new} may reverse
the sign of $J_\sigma$ near the central crossing. The next lemma controls
this contribution at central heights that are large relative to the
boundary coefficients.

\begin{lemma}[Large central heights]
\label{lem:relative-height-paired-new}
After fixing $z_*$ sufficiently small, there are constants
$\theta_0,c_*>0$ such that the following holds.  Suppose $q_1>0$,
$0<h<h_0$, and
\begin{equation}
\label{eq:relative-height-paired-condition-new}
 \frac{|q_0|}{h}+\frac{q_1}{\sqrt h}+|q_2|\le\theta_0.
\end{equation}
Then both halves through $(0,h)$ reach the endpoint windows and
\begin{equation}
\label{eq:relative-height-local-bound-new}
 0<L_p'\le\exp(-c_*\sqrt{z_*})<1.
\end{equation}
The return is strictly contracting in $I$, and $\Delta_p'(h)<0$ if
$h\in\mathcal H_p$.
\end{lemma}

\begin{proof}
We first continue both graphs in a cone, then place their endpoints in
the windows, and finally estimate the derivative ratio. A first-contact
comparison in
\eqref{eq:physical-graph-equation-new}, first in a provisional fixed box
and then in the resulting cone, gives constants $c,C>0$ such that
\begin{equation}
\label{eq:large-height-cone-new}
 c(z-h)\le U_\sigma(z;h)\le C(z-h),\qquad
 \mathcal G_{\sigma,p}(U_\sigma,z)\ge c,
 \qquad h\le z\le z_*.
\end{equation}
Indeed, in this cone the constant, linear, quadratic, and cubic parts of
$q/z$ are bounded, respectively, by
\[
 \frac{|q_0|}{h},\qquad
 C\frac{q_1}{\sqrt h},\qquad
 C|q_2|,\qquad C\sqrt{z_*}.
\]
Choose $z_*$ and then $\theta_0$ so that these quantities fit the strict
cone margins. This proves continuation to $z=z_*$ and
\eqref{eq:large-height-cone-new}.

To place the endpoints in the prescribed windows, choose a fixed
intermediate height \(\zeta\ll z_*\), then reduce \(h_0\) so that
\(h_0<\zeta/2\). The cone gives \(U_\sigma(\zeta)=O(\zeta)\),
and the organizing boundary graph has the same bound. On
\([\zeta,z_*]\), continuous dependence and the first-variation
estimate show that the endpoint difference between these graphs is
\(O(\zeta)+o_p(1)\). The organizing variational coefficient on its
positive cone is \(O(z^{-1/2})\), whose integral is uniformly bounded;
for fixed \(\zeta\), the perturbed coefficients converge uniformly on
this interval. Choose \(\zeta\) small relative to the fixed window
widths, and then reduce the parameter neighborhood.

With continuation and the endpoint windows established, we return to
the derivative ratio. The sign margins \eqref{eq:normal-sign-margins-new} give
$J_\sigma\ge b_1-q_1/z$.
In \eqref{eq:paired-log-derivative-new}, use the upper cone bound for
the negative contribution and the lower cone bound for the positive
contribution. Since
\[
 \int_h^{z_*}\frac{\dd z}{\sqrt{z-h}}=2\sqrt{z_*-h},
 \qquad
 \int_h^{z_*}\frac{\dd z}{z\sqrt{z-h}}
 \le\frac{\pi}{\sqrt h},
\]
we obtain
\begin{equation}
\label{eq:relative-height-log-bound-new}
 \log L_p'
 \le-c\sqrt{z_*-h}+C\frac{q_1}{\sqrt h}
 \le-c'\sqrt{z_*}+C\theta_0,
\end{equation}
where $h<h_0<z_*/2$. Decreasing $\theta_0$ absorbs the positive term and proves
\eqref{eq:relative-height-local-bound-new}.  Composition with
\eqref{eq:exterior-uniform-contraction-new} proves the last assertion.
\end{proof}
